% ============================================================
% SIRCA-RAG — Springer LNCS (Español, versión en castellano)
% WAIMLAp 2026 / ACSAR
% ============================================================
\documentclass[runningheads]{llncs}

\usepackage[utf8]{inputenc}
\usepackage[T1]{fontenc}
\usepackage[spanish,es-tabla,es-noshorthands]{babel}
\usepackage{amsmath,amssymb}
\usepackage{graphicx}
\usepackage{booktabs}
\usepackage{multirow}
\usepackage{array}
\usepackage{enumitem}
\usepackage{url}
\usepackage[dvipsnames]{xcolor}
\usepackage{tikz}
\usetikzlibrary{shapes.geometric,arrows.meta,positioning,calc}
\usepackage[colorlinks=true,
            linkcolor=blue!70!black,
            citecolor=blue!70!black,  
            urlcolor=blue!60!black]{hyperref}

\graphicspath{{figures/}}

\newcommand{\sirca}{\textsc{SIRCA-RAG}}

% ============================================================
\title{Corrective RAG H\'{i}brido para generaci\'{o}n
fundamentada sobre plantas medicinales Peruanas}

\titlerunning{Corrective RAG H\'{i}brido para Plantas Medicinales Peruanas}

\author{Ower Lopez-Arela\orcidID{0009-0008-1075-4275} \and
Daniel Marron-Carcausto\orcidID{0009-0008-5964-0201} \and
Antonio Arroyo-Paz\orcidID{0000-0003-1838-4527}}

\authorrunning{O. Lopez-Arela, D. Marron-Carcausto, A. Arroyo-Paz}

\institute{Escuela Profesional de Ingenier\'{i}a de Sistemas, Universidad Nacional de San Agust\'{i}n de Arequipa, Per\'{u}\\
\email{\{olopeza, dmarron, aarroyop\}@unsa.edu.pe}}

% ============================================================
\begin{document}
\sloppy
\maketitle

% ============================================================
\begin{abstract}
La vasta literatura existente sobre las plantas medicinales del Per\'{u} sufre un grave problema de fragmentaci\'{o}n. Al estar repartida entre numerosos repositorios, extraer evidencia farmacol\'{o}gica real resulta costoso. Para hacer frente a esto, proponemos \sirca{} (Sistema para la Recuperaci\'{o}n Inteligente y Respuestas Correctivas), una arquitectura biling\"{u}e dise\~{n}ada para unificar dicha evidencia mediante t\'{e}cnicas correctivas. Nuestra soluci\'{o}n fusiona b\'{u}squedas densas con l\'{e}xicas (BM25) usando \emph{Reciprocal Rank Fusion}, y luego aplica un reordenador cruzado (\emph{cross-encoder}) que termina siendo vital para proteger la exactitud de las respuestas. Un enrutador tipo CRAG gestiona la calidad final. Evaluamos la herramienta en un entorno real con 16{,}486 art\'{i}culos (cubriendo 100 especies locales), logrando una Fidelidad de $0{,}566$ y un MRR de $0{,}837$. Tambi\'{e}n demostramos emp\'{i}ricamente que la normalizaci\'{o}n por defecto arruina los enrutamientos; tuvimos que forzar una calibraci\'{o}n sigmoide absoluta para que las derivaciones hacia b\'{u}squeda externa funcionaran de forma leg\'{i}tima.

\keywords{Retrieval-Augmented Generation (RAG) \and Etnobot\'{a}nica Peruana \and Modelos de Lenguaje Grande (LLM) \and Recuperaci\'{o}n H\'{i}brida \and Procesamiento de Lenguaje Natural (NLP)}
\end{abstract}

% ============================================================
\section{Introducci\'{o}n}
\label{sec:intro}

Sabemos que el Per\'{u} posee cerca de 25{,}000 especies vegetales, y al menos entre 1{,}400 y 5{,}000 gozan de alg\'{u}n grado de evidencia farmacol\'{o}gica~\cite{bussmann2006traditional,bussmann2011toxicity}. Sin embargo, si un investigador quiere rastrear datos espec\'{i}ficos sobre \textit{Uncaria tomentosa} o \textit{Croton lechleri}, pronto se topar\'{a} con el problema central: los datos no est\'{a}n en un solo lugar. Las publicaciones se hallan regadas por todo PubMed, as\'{i} como en bases de biodiversidad locales y registros internacionales~\cite{lock2016bioactive,gonzales2012ethnobiology}. Armar una revisi\'{o}n sistem\'{a}tica en estas condiciones toma demasiado tiempo y deja vac\'{i}os.

\'{U}ltimamente, los Modelos de Lenguaje (\emph{LLMs}) han cobrado inter\'{e}s para resumir grandes vol\'{u}menes de textos. Al inyectarles contexto a trav\'{e}s de Generaci\'{o}n Aumentada por Recuperaci\'{o}n (RAG)~\cite{lewis2020retrieval}, el modelo parece ``saber'' del tema. No obstante, en la pr\'{a}ctica m\'{e}dica esto choca contra la pared: a menudo el sistema jala p\'{a}rrafos que no sirven de nada y, peor a\'{u}n, inventa conclusiones farmacol\'{o}gicas que lucen verdaderas pero carecen de respaldo~\cite{gao2024retrieval,ji2023survey}. Ni siquiera los modelos entrenados espec\'{i}ficamente en textos m\'{e}dicos se salvan de esto~\cite{lee2020biobert,gu2021domain}. Para sortear este escollo, Yan et al.~\cite{yan2024crag} publicaron Corrective RAG (CRAG). La idea detr\'{a}s de CRAG es simple pero efectiva: medir si lo que acabas de buscar sirve o no; si es malo, el modelo pide reescribir la pregunta o se va a buscar en la web.

El aporte central que presentamos aqu\'{i} es \sirca{} (Sistema para la Recuperaci\'{o}n Inteligente y Respuestas Correctivas). Lo que hicimos fue tomar ese concepto de CRAG y adaptarlo para que maneje la enorme complejidad bot\'{a}nica del Per\'{u}. Tuvimos que tomar cuatro decisiones fuertes de dise\~{n}o. Primero, decidimos no casarnos con un solo m\'{e}todo de b\'{u}squeda; combinamos vectores de \texttt{multilingual-e5-base}~\cite{wang2024e5} con la b\'{u}squeda cl\'{a}sica de palabras clave BM25~\cite{robertson2009bm25} empleando RRF~\cite{cormack2009rrf}. Segundo, metimos un \emph{cross-encoder}~\cite{nogueira2019passage} porque nos dimos cuenta que los bi-encoders b\'{a}sicos se saltan detalles t\'{e}cnicos sutiles. Tercero, implementamos un candado en el enrutamiento: si la confianza cae, bloqueamos al generador. Cuarto, forzamos al modelo (v\'{i}a CoT~\cite{wei2022chain}) a extraer los datos, verificarlos y reci\'{e}n entonces escribir la respuesta indicando qu\'{e} pasaje us\'{o}. Todo esto nos dej\'{o} con cinco contribuciones puntuales: una arquitectura QA biling\"{u}e para nuestra flora; un corpus enorme curado de 16{,}486 art\'{i}culos abarcando 100 especies; 50 preguntas validadas por nosotros mismos; una forma de medir la fidelidad combinando l\'{e}xico con sem\'{a}ntica; y finalmente, pruebas directas con m\'{u}ltiples modelos.

% ============================================================
\section{Trabajos Relacionados}
\label{sec:related}

\textbf{Ecosistema RAG.} Todo arranc\'{o} con Lewis et al.~\cite{lewis2020retrieval}, quienes demostraron que darle documentos a un modelo mejora lo que genera. A\~{n}os despu\'{e}s, Gao et al.~\cite{gao2024retrieval} hicieron un mapa de todas las variantes que hab\'{i}an salido. De todas ellas, la propuesta de Yan et al.~\cite{yan2024crag} captur\'{o} nuestra atenci\'{o}n por su enfoque correctivo: medir la calidad antes de generar. Nosotros tomamos esa ruta, pero en lugar de entrenar un evaluador desde cero, usamos directamente los \emph{logits} de un \emph{cross-encoder}. Existe tambi\'{e}n Self-RAG~\cite{asai2024selfrag}, que hace que el LLM decida cu\'{a}ndo necesita m\'{a}s datos generando tokens especiales, pero requiere modificar internamente el LLM y eso sale de nuestro presupuesto de recursos.

\textbf{B\'{u}squeda H\'{i}brida.} Hay suficiente evidencia de que mezclar enfoques densos y dispersos funciona mejor que usarlos por separado~\cite{karpukhin2020dense}. Por eso optamos por RRF~\cite{cormack2009rrf} para juntar \texttt{e5-base}~\cite{wang2024e5} con BM25~\cite{robertson2009bm25}. Lo bueno de esto es que no nos pide datos etiquetados para calibrar pesos. Despu\'{e}s, simplemente le pasamos el resultado a un \emph{cross-encoder}~\cite{nogueira2019passage}.

\textbf{C\'{o}mo Evaluamos.} Usamos el marco RAGAs~\cite{es2023ragas} porque ya es un est\'{a}ndar. Adem\'{a}s, metimos un mecanismo tipo \emph{LLM-as-judge}~\cite{zheng2023mtbench} (lo detallamos en la Sec.~\ref{sec:llm_judge}) para tener una especie de auditor\'{i}a cruzada independiente.

\textbf{El Tema M\'{e}dico Local.} Ya Xiong et al.~\cite{xiong2024mirage} con su proyecto MIRAGE nos hab\'{i}an mostrado que RAG le gana a los modelos crudos en temas m\'{e}dicos. Mirando a nivel Latinoam\'{e}rica, Monroy Barrios et al.~\cite{monroy2026finetuning} agarraron textos de flora peruana y ajustaron un modelo directamente; bajaron la perplejidad brutalmente, pero el conocimiento se queda atascado en las neuronas del modelo. Nosotros preferimos RAG porque si algo sale mal, podemos rastrear el documento original, algo vital si hablamos de farmacolog\'{i}a. En esa misma l\'{i}nea de arquitecturas alternativas, Collanqui et al.~\cite{collanqui2026cagrag} compararon RAG contra \emph{Cache-Augmented Generation} (CAG) para un chatbot de admisi\'{o}n universitaria: un \'{u}nico reglamento institucional precargado en el contexto del LLM, con caché sem\'{a}ntico para consultas repetidas. Concluyeron que CAG supera a su RAG (configurado deliberadamente con $k{=}1$, un solo fragmento por consulta) en ese dominio est\'{a}tico y acotado. No adoptamos esa comparaci\'{o}n num\'{e}ricamente: su corpus (un documento), dominio (administrativo) y estrategia de recuperaci\'{o}n difieren demasiado de los nuestros (16{,}486 art\'{i}culos, 100 especies) como para que una cifra aislada sea informativa; la mencionamos aqu\'{i} porque plantea una alternativa arquitect\'{o}nica leg\'{i}tima cuando el dominio es peque\~{n}o y estable, a diferencia del nuestro. Buscamos exhaustivamente y nadie hab\'{i}a juntado un flujo CRAG con recuperaci\'{o}n h\'{i}brida para plantas curativas andinas.

% ============================================================
\section{Arquitectura del Sistema}
\label{sec:architecture}

\sirca{} ejecuta cinco pasos secuenciales por consulta, orquestados en Python con FastAPI y LangGraph (Fig.~\ref{fig:architecture}).

\begin{figure}[!ht]
  \centering
  \begin{tikzpicture}[
    font=\sffamily\scriptsize,
    >={Latex[length=2mm,width=1.6mm]},
    phase/.style={rectangle, rounded corners=3pt, draw=black!75, line width=0.7pt,
      minimum height=9mm, text width=25mm, align=center, fill=blue!10, inner sep=2pt},
    io/.style={rectangle, rounded corners=3pt, draw=black!75, line width=0.7pt,
      fill=orange!18, minimum height=9mm, text width=13mm, align=center},
    flow/.style={->, line width=0.7pt, draw=black!75},
    branch/.style={->, line width=0.5pt, draw=red!65!black, dashed},
    blab/.style={font=\scriptsize\itshape, text=red!65!black, fill=white, inner sep=1pt}
  ]
  \node[io] (query) at (0,0) {\textbf{Consulta}\\{\scriptsize EN/ES}};
  \node[phase, right=3mm of query] (cls)
    {\textbf{1. Clasificar}\\{\scriptsize factual/explor./comp.}\\{\scriptsize $\rightarrow$ fija $\alpha$}};
  \node[phase, right=3mm of cls] (hyb)
    {\textbf{2. Recup. H\'{i}brida}\\{\scriptsize e5-base+BM25}\\{\scriptsize RRF, top-30}};
  \node[phase, right=3mm of hyb] (crag)
    {\textbf{3. Rerank+CRAG}\\{\scriptsize Cross-enc. MiniLM}\\{\scriptsize eval. $s_{\sigma}{=}\sigma(s_{\max})$}};
  \draw[flow] (query)--(cls); \draw[flow] (cls)--(hyb); \draw[flow] (hyb)--(crag);
  \node[phase, below=8mm of cls.south, anchor=north] (ctx)
    {\textbf{4. Contexto}\\{\scriptsize Top-5 docs+citas}};
  \node[phase, right=3mm of ctx] (gen)
    {\textbf{5. Generaci\'{o}n Fund.}\\{\scriptsize CoT: Extr.$\rightarrow$Verif.}\\{\scriptsize $\rightarrow$Componer (DeepSeek)}};
  \node[io, right=3mm of gen] (ans) {\textbf{Respuesta}\\{\scriptsize + citas}};
  \draw[flow] (ctx)--(gen); \draw[flow] (gen)--(ans);
  \draw[flow] (crag.south)--++(0,-4mm)-|(ctx.north)
    node[pos=0.2, font=\scriptsize\bfseries, fill=white, inner sep=1pt]{aceptar ($s_{\sigma}{\geq}0{,}60$)};
  \draw[branch] (crag.north)--++(0,4mm)-|(cls.north)
    node[pos=0.5, blab, anchor=south]{refinar ($0{,}30{\leq}s_{\sigma}{<}0{,}60$) / b\'{u}squeda web ($s_{\sigma}{<}0{,}30$)};
  \end{tikzpicture}
  \caption{Pipeline de cinco fases de \sirca{}. La clasificaci\'{o}n fija $\alpha$; la recuperaci\'{o}n h\'{i}brida usa RRF; el reordenamiento con cross-encoder alimenta el enrutamiento CRAG; y la generaci\'{o}n fundamentada aplica Chain-of-Thought. Las flechas discontinuas marcan las rutas correctivas de CRAG.}
  \label{fig:architecture}
\end{figure}

El proceso completo se divide en fases conectadas entre s\'{i}. Lo primero que hace el sistema es leer la pregunta y decidir qu\'{e} tipo de consulta es. Dependiendo de los patrones verbales que encontremos, la encasillamos como \textit{factual}, \textit{exploratoria} o \textit{comparativa}. Esta clasificaci\'{o}n no es un mero adorno; altera internamente el peso $\alpha$ que le damos a la b\'{u}squeda. Si nos piden un dato duro, le damos m\'{a}s peso a la coincidencia de palabras clave v\'{i}a BM25 ($\alpha{=}0{,}3$). Si la duda es abierta, privilegiamos el lado sem\'{a}ntico ($\alpha{=}0{,}7$).

Terminada la clasificaci\'{o}n, arranca la b\'{u}squeda paralela. Por una pista corre \texttt{multilingual-e5-base}~\cite{wang2024e5}, transformando el texto en vectores de 768 dimensiones que se comparan matem\'{a}ticamente en FAISS~\cite{johnson2021billion}. Por la otra pista corre BM25~\cite{robertson2009bm25}, buscando t\'{e}rminos espec\'{i}ficos que el vector podr\'{i}a haber ignorado. Ambas listas de resultados chocan y se entrelazan mediante una simple f\'{o}rmula posicional llamada RRF~\cite{cormack2009rrf}, dej\'{a}ndonos con una canasta de 30 textos probables.

Ah\'{i} entra a tallar el \texttt{ms-marco-MiniLM-L-12-v2}~\cite{nogueira2019passage}. Este m\'{o}dulo toma los 30 textos y la pregunta original, y los lee juntos para afinar el puntaje de relevancia. Nos quedamos solo con los 10 mejores. En este punto, el sistema eval\'{u}a la nota m\'{a}xima obtenida. Pasamos esa nota por una funci\'{o}n sigmoide para aplanarla y decidir: si sobrepasa el $0{,}60$ (y el resto de documentos no es deficiente), aprobamos el lote. Si cae entre $0{,}30$ y $0{,}60$, levantamos una alerta de ``refinamiento''. Y si es menor a $0{,}30$, el sistema claudica con su base de datos local y emite la se\~{n}al para recurrir a internet. En la secci\'{o}n posterior explicamos por qu\'{e} esta decisi\'{o}n matem\'{a}tica de usar la sigmoide result\'{o} vital para el dise\~{n}o.

Finalmente, si el lote fue aprobado, tomamos los 5 mejores extractos y se los pasamos al modelo DeepSeek V4-Flash~\cite{deepseek2026v4}\footnote{Accedido a trav\'{e}s de la API oficial de DeepSeek (\url{https://api.deepseek.com}), bajo el identificador de modelo \texttt{deepseek-v4-flash}.}. No le exigimos la respuesta sin m\'{a}s; le imponemos un flujo l\'{o}gico de extraer, comparar, y reci\'{e}n redactar. Y lo m\'{a}s cr\'{i}tico: se bloquea expl\'{i}citamente su capacidad para inferir datos que no est\'{e}n presentes en esos p\'{a}rrafos.

% ============================================================
\section{Corpus y Evaluaci\'{o}n}
\label{sec:experiments}

\textbf{Sobre los datos manejados.} Recorrimos ocho repositorios grandes para recopilar esto. Revisamos PubMed, Europe PMC, Semantic Scholar y CrossRef, adem\'{a}s de bases duras como GBIF, PeruNPDB, WFO y COCONUT. Nos enfocamos en 100 plantas del sur peruano. Quitamos todos los art\'{i}culos repetidos cruzando los DOIs y nos quedamos con 16{,}486 documentos \'{u}nicos. Partimos todo ese texto en pedazos de 512 palabras (con cierto solape para no cortar ideas por la mitad). Noventa y una de esas plantas s\'{i} ten\'{i}an literatura; el resto (\textit{Adesmia spinosissima}, \textit{Berberis lutea}, entre otras flores de altura) no arrojaron pr\'{a}cticamente nada \'{u}til. Esto lo documentamos a prop\'{o}sito para testear c\'{o}mo falla el sistema m\'{a}s adelante. Hemos publicado todo este repositorio en l\'{i}nea para quien quiera revisarlo.

\textbf{Acerca de los experimentos.} Creamos nosotros mismos un set de 50 preguntas en espa\~{n}ol e ingl\'{e}s (mezclando preguntas de concepto, comparaciones y datos duros) sobre 12 especies y validamos cu\'{a}l deber\'{i}a ser la respuesta ideal. Para no quedarnos solo con eso, armamos la evaluaci\'{o}n en tres frentes: las 50 preguntas base; un set extendido que abarca 30 especies adicionales para ver si el sistema no se marea; y un paquete de preguntas que dise\~{n}amos con malicia para obligar al sistema a fracasar y activar el rescate web.

En cuanto a c\'{o}mo medimos el \'{e}xito, nos guiamos por par\'{a}metros est\'{a}ndar. Del lado de recuperar los textos usamos Context Precision, Recall, MRR y NDCG. Del lado de la redacci\'{o}n, usamos BERTScore F1~\cite{zhang2020bertscore}. Y para estar tranquilos de que el modelo no est\'{a} alucinando datos m\'{e}dicos, ideamos una m\'{e}trica de ``Fidelidad'' que mezcla evaluaci\'{o}n sem\'{a}ntica con coincidencias literales de texto. Tambi\'{e}n corrimos simulaciones quitando piezas del sistema paso a paso (lo que llamamos ablaci\'{o}n) para probar qu\'{e} pasaba.

% ============================================================
\section{Resultados}
\label{sec:results}

Sobre el \emph{benchmark} de 50 consultas y 12 especies humano-verificadas, \sirca{} obtiene un Context Recall de $0{,}492$, MRR de $0{,}837$ y Fidelidad de $0{,}566$ (Tabla~\ref{tab:full_results}, columna \texttt{full}). El MRR indica que, en la gran mayor\'{i}a de las consultas, el documento m\'{a}s relevante queda en la primera o segunda posici\'{o}n del ranking. No comparamos estos valores num\'{e}ricamente contra antecedentes publicados: el \'{u}nico trabajo af\'{i}n que identificamos en espa\~{n}ol (Collanqui et al.~\cite{collanqui2026cagrag}, Sec.~\ref{sec:related}) eval\'{u}a un dominio, corpus y estrategia de recuperaci\'{o}n demasiado distintos —un \'{u}nico reglamento institucional con recuperaci\'{o}n $k{=}1$— como para que una comparaci\'{o}n cuantitativa directa sea v\'{a}lida.

Las m\'{e}tricas de generaci\'{o}n tambi\'{e}n son s\'{o}lidas: BERTScore F1 de $0{,}841$ (\texttt{roberta-large}), Similitud Sem\'{a}ntica de $0{,}852$ y Relevancia de Respuesta de $0{,}998$. La excepci\'{o}n es la Cobertura de Entidades ($0{,}402$), que queda por debajo del resto. Inspeccionamos manualmente qu\'{e} ocurr\'{i}a con las consultas que registraron notas m\'{a}s bajas. Descubrimos que el paso final del generador es en extremo selectivo: suele mencionar un \'{u}nico compuesto o planta por afirmaci\'{o}n (aquel con mayor soporte textual) e ignora entidades secundarias que s\'{i} figuran en el p\'{a}rrafo de referencia. Como resultado, las m\'{e}tricas automatizadas que exigen exhaustividad de entidades lo penalizan, a pesar de que la respuesta sea t\'{e}cnicamente correcta. La Fidelidad de $0{,}566$, por su parte, refleja el efecto del componente l\'{e}xico del 35\%: ese factor penaliza a prop\'{o}sito la par\'{a}frasis de f\'{o}rmulas qu\'{i}micas y nombres taxon\'{o}micos para reducir el riesgo de alucinaci\'{o}n farmacol\'{o}gica.

\begin{table}[!ht]
\centering
\caption{M\'{e}tricas de recuperaci\'{o}n (pool top-30, RRF h\'{i}brido + cross-encoder, misma metodolog\'{i}a que la Tabla~\ref{tab:full_results}) sobre el subconjunto inicial ($n{=}12$ especies humano-verificadas) y una extensi\'{o}n de 30 especies adicionales muestreadas del cat\'{a}logo (42 especies \'{u}nicas en total).}
\label{tab:extended_retrieval}
\resizebox{0.9\linewidth}{!}{%
\begin{tabular}{l c c c c c c}
\toprule
\textbf{Subconjunto} & $n_{\text{spp}}$ & \textbf{MRR} & \textbf{NDCG@5} & \textbf{NDCG@10} & \textbf{C.\ Precision} & \textbf{C.\ Recall} \\
\midrule
Inicial humano-verificado & 12 & $0{,}837$ & $0{,}757$ & $0{,}750$ & $0{,}415$ & $0{,}492$ \\
Extensi\'{o}n (30 especies nuevas) & 30 & $0{,}870$ & $0{,}816$ & $0{,}841$ & $0{,}463$ & $0{,}535$ \\
\bottomrule
\end{tabular}%
}
\end{table}

\textbf{Generalizaci\'{o}n de la recuperaci\'{o}n.} La Tabla~\ref{tab:extended_retrieval} muestra que la calidad de la recuperaci\'{o}n se sostiene —y de hecho mejora ligeramente— al ampliar el marco de evaluaci\'{o}n a 30 especies no vistas del corpus: sobre ese conjunto el MRR sube de $0{,}837$ a $0{,}870$, el NDCG@10 sube de $0{,}750$ a $0{,}841$ y el Context Recall tambi\'{e}n mejora ($0{,}492 \rightarrow 0{,}535$). No interpretamos esto como evidencia de que las 30 especies nuevas sean intr\'{i}nsecamente m\'{a}s f\'{a}ciles, sino como confirmaci\'{o}n de que el motor de recuperaci\'{o}n no est\'{a} sobreajustado al subconjunto de 12 especies humano-verificadas: generaliza a datos in\'{e}ditos sin degradaci\'{o}n.

\subsection{Robustez ante M\'{u}ltiples Modelos}
\label{sec:multi_llm}

Surgi\'{o} la interrogante de si lo observado era un artefacto exclusivo del modelo DeepSeek. Para descartarlo, alteramos el motor y configuramos a Cerebras Gemma-4-31B~\cite{googledeepmind2026gemma4} para responder las mismas 50 preguntas. Los n\'{u}meros de b\'{u}squeda no variaron por ser previos en la cadena, pero la redacci\'{o}n s\'{i} lo hizo de manera medible (Tabla~\ref{tab:multi_llm}).

\begin{table}[!ht]
\centering
\caption{Comparaci\'{o}n emp\'{i}rica del motor generador (DeepSeek V4-Flash vs. Cerebras Gemma-4-31B) manteniendo fija la canalizaci\'{o}n \sirca{}. Pruebas $t$ pareadas.}
\label{tab:multi_llm}
\resizebox{0.9\linewidth}{!}{%
\begin{tabular}{l c c c l}
\toprule
\textbf{M\'{e}trica} & \textbf{DeepSeek-V4} & \textbf{Gemma-4-31B} & $\Delta$ & \textbf{Significancia ($t$-pareada)} \\
\midrule
BERTScore F1          & $0{,}841$ & $0{,}830$ & $-0{,}011$ & \textbf{Significativo} ($p{=}0{,}005$) \\
Similitud Sem\'{a}ntica & $0{,}852$ & $0{,}789$ & $-0{,}063$ & No significativo ($p{=}0{,}18$) \\
Cobertura de Entidades& $0{,}402$ & $0{,}386$ & $-0{,}016$ & No significativo ($p{=}0{,}46$) \\
Relevancia            & $0{,}998$ & $0{,}988$ & $-0{,}010$ & \textbf{Significativo} ($p{=}0{,}042$) \\
\midrule
\textbf{Fidelidad} (h\'{i}brida) & $0{,}517$ & $\mathbf{0{,}609}$ & $\mathbf{+0{,}092}$ & \textbf{Significativo} ($p{<}0{,}001$) \\
\bottomrule
\end{tabular}%
}
\end{table}

Los resultados muestran un patr\'{o}n m\'{a}s matizado de lo que anticip\'{a}bamos: DeepSeek V4-Flash iguala o supera a Gemma-4-31B en cuatro de las cinco m\'{e}tricas de generaci\'{o}n —incluyendo BERTScore F1 ($p{=}0{,}005$) y Relevancia ($p{=}0{,}042$), ambas significativas a favor de DeepSeek—, mientras que Similitud Sem\'{a}ntica y Cobertura de Entidades no muestran diferencias significativas. La \'{u}nica m\'{e}trica donde Gemma es netamente superior es la Fidelidad ($p{<}0{,}001$). Atribuimos esta \'{u}nica ventaja al estilo compositivo m\'{a}s literal de Gemma, que tiende a reproducir nombres cient\'{i}ficos y frases del contexto casi verbatim —lo cual beneficia espec\'{i}ficamente al componente l\'{e}xico (35\,\%) de nuestra Fidelidad h\'{i}brida—, mientras que en las m\'{e}tricas de calidad general (BERTScore, Relevancia) la redacci\'{o}n m\'{a}s fluida de DeepSeek no se ve penalizada. El punto arquitect\'{o}nicamente relevante es doble: (i)~la pipeline \sirca{} tolera el reemplazo del generador sin ajustes adicionales, y (ii)~la elecci\'{o}n del generador s\'{i} importa para la Fidelidad espec\'{i}ficamente, no para la calidad de generaci\'{o}n en general. A diferencia de una corrida temprana de este mismo experimento, esta corrida no registr\'{o} respuestas vac\'{i}as en ninguno de los dos modelos ($0/50$ cada uno), tras el fix de reintento descrito en la Limitaci\'{o}n~(vii).

\subsection{Verificaci\'{o}n Cruzada}
\label{sec:llm_judge}

Quer\'{i}amos estar seguros de que nuestra m\'{e}trica de Fidelidad no nos estuviera mintiendo. As\'{i} que aplicamos la t\'{e}cnica \emph{LLM-as-judge}~\cite{zheng2023mtbench}, poniendo al propio Gemma-4 a calificar del 0 al 4 qu\'{e} tan bien hab\'{i}a respondido DeepSeek respecto a los textos entregados. Hicimos a prop\'{o}sito que el modelo evaluador fuera de otra familia para evitar sesgos amigables. El resultado cualitativo es alentador: 43 de las 50 respuestas sacaron puntaje perfecto (4/4), una obtuvo 3, cinco obtuvieron 2, y una sola recibi\'{o} cero por no aportar informaci\'{o}n \'{u}til. Ese $86\%$ de calificaciones perfectas deja, sin embargo, muy poca varianza en el juez para una correlaci\'{o}n formal: el coeficiente de Pearson entre el puntaje del juez y nuestra m\'{e}trica algor\'{i}tmica result\'{o} bajo y no significativo ($r{=}0{,}056$, $p{=}0{,}70$), esperable cuando una de las dos variables est\'{a} concentrada casi por completo en su valor m\'{a}ximo (efecto techo).

Lo interesante aqu\'{i} es la brecha cruda de los n\'{u}meros. El juez promediaba notas altas ($0{,}925$ normalizado), pero nuestra m\'{e}trica era mucho m\'{a}s dura ($0{,}517$). Esto ocurri\'{o} porque dise\~{n}amos nuestra f\'{o}rmula para ser implacable con los parafraseos m\'{e}dicos: preferimos que el sistema nos reporte una baja nota antes de dar por v\'{a}lida una frase que suene ``parecida'' pero que podr\'{i}a contener un error cl\'{i}nico. Reconocemos que ning\'{u}n juez artificial va a reemplazar el rigor toxicol\'{o}gico de un farmac\'{e}utico humano; el ejercicio funciona mejor como filtro cualitativo grueso (detectar respuestas sin sustento) que como validaci\'{o}n cuantitativa estricta de la Fidelidad, dado el efecto techo observado.

\begin{table}[!ht]
\centering
\caption{Ablation de seis configuraciones (50 consultas). Recuperaci\'{o}n: metodolog\'{i}a basada en el agente con el evaluador CRAG calibrado por sigmoide (script \texttt{run\_perquery\_agent\_ablation.py}). Fidelidad: reportada para \texttt{full} y \texttt{no\_reranker} con generaci\'{o}n DeepSeek V4-Flash (script \texttt{run\_n5\_fidelity\_wilcoxon.py}); ``n/d'' indica configuraciones cuya re-ejecuci\'{o}n LLM por-consulta queda como trabajo futuro (Limitaci\'{o}n iii). Negrita = mejor; subrayado = peor por columna. $^\dagger$\texttt{dense\_only} coincide con \texttt{full} en recuperaci\'{o}n porque tras el reordenamiento el Top-10 es el mismo en las 50 consultas (invariante a $\alpha$ cuando el cross-encoder domina); \texttt{no\_classifier} \textbf{no} comparte esa invariancia (ver Secci\'{o}n~\ref{sec:wilcoxon}).}
\label{tab:full_results}
\footnotesize
\begin{tabular}{l c c c c c c}
\toprule
\textbf{Config.} & \textbf{C.Prec.} & \textbf{C.Recall} & \textbf{MRR} & \textbf{NDCG@10} & \textbf{Fidel.} & $\boldsymbol{\Delta}$\textbf{Fidel.} \\
\midrule+
\texttt{full}           & 0{,}411 & 0{,}482 & 0{,}825 & 0{,}741 & \textbf{0{,}566} & --- \\
\texttt{dense\_only}$^\dagger$ & 0{,}411 & 0{,}482 & 0{,}825 & 0{,}741 & n/d & n/d \\
\texttt{sparse\_only}   & \underline{0{,}402} & \underline{0{,}466} & \underline{0{,}752} & \underline{0{,}666} & n/d & n/d \\
\texttt{no\_reranker}   & 0{,}418 & 0{,}476 & \textbf{0{,}869} & 0{,}751 & \underline{0{,}465} & $-17{,}8\%$ \\
\texttt{no\_crag}       & 0{,}431 & 0{,}511 & 0{,}851 & 0{,}772 & n/d & n/d \\
\texttt{no\_classifier} & \textbf{0{,}479} & \textbf{0{,}558} & 0{,}862 & \textbf{0{,}833} & n/d & n/d \\
\bottomrule
\end{tabular}
\end{table}

\subsection{Hallazgos al Desarmar el Sistema}

Al quitar piezas una por una (Tabla~\ref{tab:full_results}), saltaron a la vista cuatro verdades inc\'{o}modas pero \'{u}tiles. Primero, el m\'{o}dulo de reordenamiento es literalmente el escudo del sistema. Cuando lo apagamos, la Fidelidad se desplom\'{o} un $17{,}8\%$ relativo (de $0{,}566$ a $0{,}465$), diferencia estad\'{i}sticamente significativa por Wilcoxon pareado ($p{=}0{,}00023$; Sec.~\ref{sec:wilcoxon}). Resulta que la fusi\'{o}n inicial s\'{i} trae textos relevantes a la parte alta de la lista, pero vienen acompa\~{n}ados de ``falsos amigos'' que terminan enga\~{n}ando al generador si el reordenador no los filtra a tiempo. Segundo, la b\'{u}squeda por palabras clave sola (\texttt{sparse\_only}) result\'{o} ser la configuraci\'{o}n m\'{a}s d\'{e}bil en las cuatro m\'{e}tricas de recuperaci\'{o}n (C.Prec. $0{,}402$, C.Recall $0{,}466$, MRR $0{,}752$, NDCG@10 $0{,}666$, todas el peor valor por columna): sin el componente denso, BM25 pierde los sin\'{o}nimos y las variantes de redacci\'{o}n de las 50 consultas. Tercero, y m\'{a}s inc\'{o}modo todav\'{i}a: desactivar el clasificador de preguntas (\texttt{no\_classifier}) \emph{mejora} la recuperaci\'{o}n de forma significativa frente a \texttt{full} —C.Prec. $+16{,}5\%$ ($p{=}0{,}007$), C.Recall $+15{,}8\%$ ($p{<}0{,}001$) y NDCG@10 $+12{,}4\%$ ($p{=}0{,}003$)—, lo que sugiere que el $\alpha$ fijo por categor\'{i}a ($0{,}3$/$0{,}7$/$0{,}5$) que asigna el clasificador est\'{a} peor calibrado que un compromiso fijo de $\alpha{=}0{,}6$ para este benchmark; discutimos esta tensi\'{o}n en la Secci\'{o}n~\ref{sec:wilcoxon}. Cuarto, \texttt{dense\_only} s\'{i} produce exactamente el mismo Top-10 que \texttt{full} en las 50 consultas ($n_{\neq}{=}0$): una vez que el cross-encoder reordena, el sesgo inicial hacia la b\'{u}squeda densa pura resulta indiferente. Mirando el reloj, la b\'{u}squeda nos toma apenas $\sim$1\,810\,ms procesar en CPU, y la redacci\'{o}n tarda entre 10 y 30 segundos usando el servidor de DeepSeek.

\subsection{Significancia Estad\'{i}stica de los Deltas del Ablation}
\label{sec:wilcoxon}

Las diferencias entre configuraciones en la Tabla~\ref{tab:full_results} son estimadores puntuales y podr\'{i}an deberse al azar. Para descartarlo, corrimos pruebas de Wilcoxon signed-rank pareadas sobre los valores per-consulta (50 consultas, test bilateral, dos colas) entre cada variante del ablation y la configuraci\'{o}n \texttt{full}. La Tabla~\ref{tab:wilcoxon} sintetiza los $p$-valores para Context Precision, Context Recall, MRR y NDCG@10.

\begin{table}[!ht]
\centering
\caption{$p$-valores Wilcoxon signed-rank (bilateral) contra \texttt{full}, calculados sobre m\'{e}tricas de recuperaci\'{o}n per-consulta con la misma metodolog\'{i}a basada en el agente que la Tabla~\ref{tab:full_results}. La prueba de significancia sobre Fidelidad se reporta separadamente para la comparaci\'{o}n central \texttt{full} vs \texttt{no\_reranker} (dos p\'{a}rrafos abajo, $p{=}0{,}00023$; script \texttt{run\_n5\_fidelity\_wilcoxon.py}); extenderla a las seis configuraciones exige re-ejecutar el LLM por cada una y queda como trabajo futuro. ``---'' indica configuraciones cuyo Top-10 recuperado es id\'{e}ntico a \texttt{full} ($n_{\neq}{=}0$ en las cuatro m\'{e}tricas). $n_{\neq}$ = n\'{u}mero de consultas con diferencia no nula en la columna C.Recall (cada m\'{e}trica tiene su propio $n_{\neq}$; se reporta el de C.Recall como referencia).}
\label{tab:wilcoxon}
\footnotesize
\begin{tabular}{l c c c c c}
\toprule
\textbf{Config.} & \textbf{$n_{\neq}$} & \textbf{C.Prec.} & \textbf{C.Recall} & \textbf{MRR} & \textbf{NDCG@10} \\
\midrule
\texttt{dense\_only}    &  0 & --- & --- & --- & --- \\
\texttt{sparse\_only}   & 40 & 0{,}450 & 0{,}523 & 0{,}241 & 0{,}100 \\
\texttt{no\_reranker}   & 24 & 0{,}637 & 0{,}786 & 0{,}345 & 0{,}695 \\
\texttt{no\_crag}       &  5 & 0{,}043 & 0{,}043 & 0{,}068 & 0{,}043 \\
\texttt{no\_classifier} & 20 & \textbf{0{,}007} & \textbf{0{,}0004} & 0{,}307 & \textbf{0{,}003} \\
\bottomrule
\end{tabular}
\end{table}

El cuadro estad\'{i}stico nos obliga a ir m\'{a}s despacio. \texttt{sparse\_only} no solo tiene el Recall agregado m\'{a}s bajo ($0{,}466$ frente a $0{,}482$ de \texttt{full}); la diferencia tampoco alcanza significancia a nivel per-consulta ($p{=}0{,}523$ en C.Recall, $n_{\neq}{=}40$): BM25 sin el componente denso pierde en unas consultas y gana en otras, y el neto no se distingue del azar pese a ser consistentemente el peor promedio. \texttt{dense\_only} s\'{i} produce exactamente el mismo Top-10 que \texttt{full} ($n_{\neq}{=}0$, marcada con ``---''): el reordenador cross-encoder vuelve la recuperaci\'{o}n final indiferente a partir de qu\'{e} $\alpha$ arranca la fusi\'{o}n. El hallazgo que rompe la intuici\'{o}n es \texttt{no\_classifier}: a diferencia de \texttt{dense\_only}, aqu\'{i} s\'{i} hay una diferencia real y \emph{significativa} contra \texttt{full} en tres de cuatro m\'{e}tricas (C.Prec. $p{=}0{,}007$, C.Recall $p{<}0{,}001$, NDCG@10 $p{=}0{,}003$; MRR no significativo, $p{=}0{,}307$), y en la direcci\'{o}n contraria a la esperada: usar un $\alpha$ fijo de $0{,}6$ para todas las consultas recupera mejor que el $\alpha$ por categor\'{i}a que asigna el clasificador. Esto no invalida la utilidad del clasificador para otros fines (p.\ ej.\ explicabilidad de la decisi\'{o}n de enrutamiento), pero s\'{i} indica que su calibraci\'{o}n actual de $\alpha$ ($0{,}3$/$0{,}7$/$0{,}5$ por categor\'{i}a) no est\'{a} optimizada para este benchmark y merece un barrido de hiperpar\'{a}metros en trabajo futuro. \texttt{no\_crag} s\'{i} difiere de \texttt{full}, aunque en solo $5$ de $50$ consultas seg\'{u}n C.Recall ($p{=}0{,}043$, significativo al $95\%$): al deshabilitar CRAG esas consultas dejan de disparar \textsc{refinar}/\textsc{b\'{u}squeda\_web} y el contexto final que recibe el generador cambia; en el resto de consultas CRAG acepta el mismo Top-10 en ambas configuraciones. La lectura de conjunto es que, al nivel de recuperaci\'{o}n, el ablation s\'{i} revela diferencias reales y en algunos casos contraintuitivas —no todas las configuraciones son estad\'{i}sticamente equivalentes—, pero el impacto m\'{a}s consistente y arquitect\'{o}nicamente central del ablation sigue viviendo en la Fidelidad, donde quitar el reordenador cuesta un $-17{,}8\%$ relativo (Tabla~\ref{tab:full_results}), una m\'{e}trica de generaci\'{o}n que las pruebas de recuperaci\'{o}n no alcanzan a capturar. Para verificar que esa ca\'{i}da no es ruido, aplicamos una prueba de Wilcoxon signed-rank pareada sobre la Fidelidad per-consulta entre \texttt{full} y \texttt{no\_reranker} (con generaci\'{o}n DeepSeek): la diferencia es estad\'{i}sticamente significativa ($p{=}0{,}00023$ a dos colas; $p{=}0{,}00012$ a una cola, con las 50/50 consultas registrando una diferencia no nula). Es decir, a diferencia de algunas m\'{e}tricas de recuperaci\'{o}n, la contribuci\'{o}n del reordenador a la Fidelidad —la afirmaci\'{o}n central del art\'{i}culo— s\'{i} supera la prueba de significancia de forma robusta.

\textbf{Nota sobre corridas y valor absoluto de la Fidelidad.} La Fidelidad de DeepSeek en la comparaci\'{o}n cross-LLM (Tabla~\ref{tab:multi_llm}, $0{,}517$) difiere de la Fidelidad de \texttt{full} en la Tabla~\ref{tab:full_results} ($0{,}566$) porque provienen de corridas distintas de la API comercial: la Tabla~\ref{tab:multi_llm} usa la ejecuci\'{o}n pareada con Gemma-4, mientras que la Tabla~\ref{tab:full_results} emplea una ejecuci\'{o}n dedicada del ablation. El resultado \'{u}nico reproducible entre corridas —y el que sustenta la conclusi\'{o}n del art\'{i}culo— es la \emph{direcci\'{o}n} y \emph{significancia} de las diferencias, no el valor absoluto de las m\'{e}tricas.

\subsection{El Enga\~{n}o de la Normalizaci\'{o}n}
\label{sec:crag_validation}

Tropezamos con un obst\'{a}culo t\'{e}cnico enorme a medio camino. La herramienta de evaluaci\'{o}n que est\'{a}bamos usando ten\'{i}a la costumbre de estirar los puntajes de cada bloque de texto de $0$ a $1$. Esto significa que, incluso si todos los textos encontrados eran basura, el "menos malo" recib\'{i}a un $1{,}0$ perfecto. Consecuentemente, el sistema daba luz verde a todo y jam\'{a}s activaba las b\'{u}squedas externas de auxilio. Tuvimos que meter mano en la arquitectura y calibrar los n\'{u}meros en crudo usando una curva sigmoide absoluta. Fijamos los m\'{a}rgenes en $0{,}60$ para aprobar y $0{,}30$ para dudar. Eso solucion\'{o} el problema de ra\'{i}z.

Para probar que esto realmente funcionaba, le tiramos al sistema un arsenal de 27 preguntas trampa. Algunas eran sobre plantas que sab\'{i}amos que no estaban en la base; otras eran absurdos inform\'{a}ticos como "Vue.js", y otras eran simplemente palabras inconexas. La Tabla~\ref{tab:crag_stress} muestra c\'{o}mo se defendi\'{o} el algoritmo.

\begin{table}[!ht]
\centering
\caption{Decisiones del sistema al procesar 27 consultas de estr\'{e}s. Familia A: Taxones sin datos indexados; Familia B: Tem\'{a}ticas ajenas al dominio bot\'{a}nico; Familia C: Cadenas l\'{e}xicas degradadas.}
\label{tab:crag_stress}
\footnotesize
\begin{tabular}{l c c c c}
\toprule
\textbf{Familia} & $n$ & \textbf{\textsc{aceptar}} & \textbf{\textsc{refinar}} & \textbf{\textsc{b\'{u}squeda\_web}} \\
\midrule
A. Especies sin datos & 9 & 1 & 3 & 5 \\
B. Fuera de dominio   & 10 & 0 & 0 & \textbf{10} \\
C. Confusas           &  8 & 4 & 0 & 4 \\
\midrule
Total (\emph{probes})            & 27 & 5 & 3 & 19 \\
Total (\emph{benchmark} en dominio) & 50 & 35 & 5 & 10 \\
\bottomrule
\end{tabular}
\end{table}

El comportamiento observado se aline\'{o} con las expectativas te\'{o}ricas. Frente a tem\'{a}ticas fuera de lugar, el sistema redirigi\'{o} el $100\%$ de las veces hacia la b\'{u}squeda externa. Al consultar por plantas ausentes en el corpus, en $8$ de $9$ intentos solicit\'{o} reformular la premisa o acudi\'{o} a internet. Hubo un \'{u}nico bypass con \textit{Senecio canescens}, ocasionado porque el motor recuper\'{o} datos de taxones filogen\'{e}ticamente afines. Sobre el propio \emph{benchmark} en dominio (50 consultas), la calibraci\'{o}n absoluta no siempre da \textsc{aceptar}: $35/50$ enrutan a \textsc{aceptar}, $5/50$ a \textsc{refinar} y $10/50$ a \textsc{b\'{u}squeda\_web}. Interpretamos estas $15$ activaciones correctivas como comportamiento correcto y no como falsos negativos: se concentran desproporcionadamente en consultas \textit{exploratorias} y \textit{comparativas} ($12/15$, $80\%$, frente a $51\%$ en las consultas aceptadas), categor\'{i}as que por dise\~{n}o piden s\'{i}ntesis a trav\'{e}s de m\'{u}ltiples fuentes en lugar de un dato puntual, y por tanto son m\'{a}s propensas a una recuperaci\'{o}n gen\'{u}inamente ambigua. Preferimos que el sistema declare incertidumbre en esos casos antes que fabricar una respuesta. En s\'{i}ntesis, demostramos de forma emp\'{i}rica que el mecanismo de redirecci\'{o}n opera con \'{e}xito y que, al reemplazar la normalizaci\'{o}n intra-batch por la calibraci\'{o}n sigmoide absoluta, la ruta correctiva se activa tanto en entradas fuera de distribuci\'{o}n como en consultas en dominio de baja confianza —el bug de la normalizaci\'{o}n previa hab\'{i}a suprimido este segundo caso, indistinguible del comportamiento nominal.

% ============================================================
\section{Discusi\'{o}n}
\label{sec:discussion}

A la hora de procesar consultas de herbolaria, nos enfrentamos a estructuras heterog\'{e}neas. Un usuario puede solicitar "nombres comunes" y concatenar un t\'{e}rmino acad\'{e}mico riguroso como \textit{Uncaria tomentosa}. Como documentamos, la b\'{u}squeda l\'{e}xica es insuperable para capturar nombres latinos, pero el m\'{e}todo h\'{i}brido se justifica porque el \emph{cross-encoder} interviene para depurar aquellos resultados que carecen de cohesi\'{o}n sem\'{a}ntica. Adicionalmente, gestionar m\'{u}ltiples idiomas de forma nativa mediante el espacio vectorial mitig\'{o} barreras computacionales considerables.

Si volteamos a ver otros esfuerzos, como el equipo de Monroy Barrios et al.~\cite{monroy2026finetuning}, ellos inyectaron a la fuerza textos andinos dentro del cerebro de Llama. Eso da respuestas flu\'{i}das, pero si un m\'{e}dico te pregunta de d\'{o}nde sacaste esa dosis, no hay forma de demostrarlo. Nuestro camino con RAG mantiene los datos afuera, auditables y con un n\'{u}mero de referencia rastreable. Creemos que el futuro obvio ser\'{a} fusionar ambos mundos: ajustar un modelo y, encima, ponerle un RAG de protecci\'{o}n.

\textbf{Disponibilidad y limitaciones.} El c\'{o}digo completo se encuentra publicado en \url{https://github.com/yughiyami/rag_medicinal_plants}; existe adem\'{a}s una demo interactiva en \url{https://rag.scn.quest} donde se pueden observar en tiempo real los fragmentos recuperados con sus DOI/PMID y la decisi\'{o}n de enrutamiento CRAG para cada consulta. Listamos a continuaci\'{o}n las limitaciones actuales con el fin de que los lectores puedan dimensionar correctamente el alcance de los resultados: (i)~el \emph{benchmark} humano-verificado abarca 12 de las 100 especies del corpus; la recuperaci\'{o}n fue verificada sobre 42 especies en total (Sec.~\ref{sec:results}) y el enrutamiento sobre las 9 especies sin datos en \'{i}ndice (Sec.~\ref{sec:crag_validation}), pero la evaluaci\'{o}n completa de calidad de respuesta sobre las 58 especies restantes exige construir respuestas de referencia adicionales; (ii)~las respuestas de referencia actuales fueron elaboradas por los propios autores; aunque la m\'{e}trica de Fidelidad fue contrastada con un LLM-juez independiente (Sec.~\ref{sec:llm_judge}), la validaci\'{o}n por un especialista externo en farmacognosia o etnobot\'{a}nica a\'{u}n est\'{a} pendiente, siendo ese el primer requisito antes de cualquier uso cl\'{i}nico o productivo; (iii)~probamos la significancia sobre Fidelidad para la comparaci\'{o}n central (\texttt{full} vs \texttt{no\_reranker}, $p{=}0{,}00023$; Sec.~\ref{sec:wilcoxon}); extender el test de Fidelidad a las seis configuraciones del ablation requiere re-ejecutar el LLM por cada una y queda como trabajo futuro; (iv)~la Fidelidad h\'{i}brida (65\% sem\'{a}ntico + 35\% l\'{e}xico) fue contrastada contra un LLM-juez en la Sec.~\ref{sec:llm_judge}; el $86\%$ de calificaciones perfectas del juez sirve como filtro cualitativo grueso, pero la correlaci\'{o}n formal no result\'{o} significativa por el efecto techo resultante ($r{=}0{,}056$, $p{=}0{,}70$), as\'{i} que esta validaci\'{o}n debe leerse como indicativa, no como prueba cuantitativa; la correlaci\'{o}n con juicios humanos expertos sigue abierta como pregunta de investigaci\'{o}n; (v)~la robustez ante distintos LLM fue verificada con un segundo generador (Cerebras Gemma-4-31B); extender esa verificaci\'{o}n a otras familias (Llama, Qwen, Mistral) fortalecer\'{i}a la afirmaci\'{o}n de generalidad; (vi)~los umbrales de calibraci\'{o}n ($\sigma_{\mathrm{acc}}{=}0{,}60$, $\sigma_{\mathrm{ref}}{=}0{,}30$) provienen de la curva sigmoide del cross-encoder y se validaron sobre 27 \emph{probes}; un barrido sistem\'{a}tico sobre un conjunto OOD m\'{a}s amplio permitir\'{i}a afinar el punto de operaci\'{o}n; (vii)~una corrida temprana de la evaluaci\'{o}n cross-LLM (Sec.~\ref{sec:multi_llm}) mostr\'{o} al generador DeepSeek V4-Flash devolviendo respuestas vac\'{i}as o truncadas en $8/50$ consultas ($16\%$); diagnosticamos la causa (ausencia de reintento ante contenido vac\'{i}o en el cliente HTTP) y la corregimos con reintento y backoff exponencial, lo que elimin\'{o} el problema por completo ($0/50$ en la corrida reportada en este art\'{i}culo); dejamos constancia del hallazgo original porque puede resurgir con otros proveedores de API que fallen de forma similar; (viii)~los valores absolutos de las m\'{e}tricas de generaci\'{o}n dependen del punto en el tiempo en que se accedi\'{o} a la API comercial de DeepSeek, por lo que la reproducci\'{o}n exacta puede diferir de los valores aqu\'{i} reportados; las conclusiones de direcci\'{o}n y significancia s\'{i} son estables entre corridas.

% ============================================================
\section{Conclusiones}
\label{sec:conclusion}

A lo largo de este manuscrito delineamos la arquitectura de \sirca{}. Integramos m\'{e}todos de b\'{u}squeda paralela, incorporamos un reordenador estricto, establecimos barreras de calibraci\'{o}n matem\'{a}tica absoluta y forzamos al modelo a emplear un razonamiento secuencial antes de sintetizar salidas bot\'{a}nicas biling\"{u}es. Tras someter el sistema al acervo de 100 especies curativas peruanas, obtuvimos un Context Recall de $0{,}492$, MRR de $0{,}837$ y Fidelidad de $0{,}566$ sobre el \emph{benchmark} de 50 consultas. Constatamos emp\'{i}ricamente que el reordenador constituye el eslab\'{o}n m\'{a}s cr\'{i}tico de la cadena de confianza; su exclusi\'{o}n reduce la Fidelidad de forma estad\'{i}sticamente significativa ($-17{,}8\%$ relativo, $p{=}0{,}00023$). Las pruebas de validaci\'{o}n cruzada con un segundo generador (Cerebras Gemma-4-31B) confirmaron que la arquitectura tolera el reemplazo del LLM sin cambios de dise\~{n}o, aunque tres de las cinco m\'{e}tricas de generaci\'{o}n —Fidelidad, BERTScore F1 y Relevancia— s\'{i} var\'{i}an de forma estad\'{i}sticamente significativa seg\'{u}n el generador elegido. Observamos, de primera mano, c\'{o}mo los comportamientos predeterminados del software est\'{a}ndar comprometen las tuber\'{i}as correctivas si no se interviene directamente en la mec\'{a}nica de la normalizaci\'{o}n emp\'{i}rica. Como corolario, la trayectoria futura es evidente: resulta imperativo estructurar paneles con bi\'{o}logos para auditar las respuestas bajo un estricto escrutinio profesional~\cite{heinrich2020ethnopharmacology}, ampliar el alcance de nuestras pruebas geogr\'{a}ficas y, en una dimensi\'{o}n m\'{a}s ambiciosa, ejecutar modificaciones param\'{e}tricas sobre los pesos del LLM para internalizar la ontolog\'{i}a bot\'{a}nica previo al an\'{a}lisis documental.

% ============================================================

% ============================================================
\bibliographystyle{splncs04}
\bibliography{references}

\end{document}
